% TOP_PROBLEM: {"id":30007601,"problem_number":"LOCAL-30007601","title":"AI trading and market efficiency","category_id":21,"category":{"id":21,"name":"economics","display_name":"Economics","description":"Open economic theory statements, published research agendas, and proposed scoped specifications; question provenance and historical priority are qualified per record.","slug":"economics","order_index":21,"created_at":"2026-10-08T00:00:00Z"},"status":"research_question","external_url":null,"legacy_ids":[],"published":false,"statement":"In the explicitly specified setting: When does widespread AI trading improve information aggregation, and when does it create collusion, correlated errors or instability? The mathematical statement above fixes the target, assumptions and scope of this catalogue formulation.","economics":{"original_id":90,"field":"Asset pricing and financial markets","kind":"empirical","formalization_basis":"adapted_research_question","source_status":"research_agenda","priority_status":"earliest_found","priority_note":"This is the earliest explicit statement verified in this audit; first priority for the full catalogue formulation is not established.","earliest_verified_reference":{"authors":["Bank of England"],"title":"Bank of England Agenda for Research, 2025–2028","year":2026,"url":"https://www.bankofengland.co.uk/research/bank-of-england-agenda-for-research","doi":null,"locator":"The Financial System / Technological innovation, paragraphs 3–4; Prudential Architecture / Resilience, last paragraph","evidence_summary":"Requests algorithmic-market effects and AI resilience frameworks."},"current_reference":{"authors":["Bank of England"],"title":"Bank of England Agenda for Research, 2025–2028","year":2026,"url":"https://www.bankofengland.co.uk/research/bank-of-england-agenda-for-research","doi":null,"locator":"The Financial System / Technological innovation, paragraphs 3–4; Prudential Architecture / Resilience, last paragraph","evidence_summary":"Requests algorithmic-market effects and AI resilience frameworks."},"formalization_note":"Adapts an explicitly verified research-agenda passage. Mathematical objects, interventions, objectives and scope are proposed here; existing model-specific solutions are not relabeled unsolved."},"statusQualification":"Published question or research agenda verified at the cited locator. The mathematical specification is adapted for this catalogue and includes additional assumptions; source publication does not establish first-ever priority or certify this exact model remains unsolved. Adapts an explicitly verified research-agenda passage. Mathematical objects, interventions, objectives and scope are proposed here; existing model-specific solutions are not relabeled unsolved.","statusReviewedAt":"2026-10-08"}
% FIRST_POSTED: 2026-10-08
% ENTRY_KIND: statement_only
% !TeX program = lualatex
\documentclass[11pt]{article}
\usepackage{fontspec}
\usepackage[a4paper,margin=25mm]{geometry}
\usepackage{amsmath,amssymb,mathtools}
\usepackage{xurl}
\usepackage[hidelinks]{hyperref}
\newcommand{\sref}[2]{\hyperlink{source:#1:#2}{[S#2]}}
\setlength{\emergencystretch}{3em}
\begin{document}
% problemId:30007601
\section*{AI trading and market efficiency}\label{30007601}
\subsection{Definitions and mathematical statement}
Fix a repeated asset market with a fundamental payoff \(V\), price \(P\), private signals, an order-processing rule and strategic traders. A learning system is specified by its training observations, objective, action space, update rule and computational restrictions. Let \(s\) denote the market share using a specified family of artificial-intelligence trading systems. Define an intervention that varies \(s\), with competitive and informational responses of other traders allowed. At share \(s\), measure price inefficiency \(E[(P-V)^2]\), effective spreads, trading volume and the expected resource loss of market disruption. Define a collusion statistic as the persistent excess of joint trading profits or spreads over a separately computed competitive benchmark under identical information and technology. The task is to identify how these outcomes change with \(s\) and characterize conditions under which learning supports tacit collusion or correlated destabilizing behavior. Vary signal precision, feedback, algorithm diversity and entry while keeping the comparison explicit. Historical backtests cannot reveal all equilibrium responses to widespread adoption; live or experimentally controlled validation must address adaptation and data leakage. Particular learning environments can exhibit inefficient outcomes without proving that all AI trading has that effect.

\par\textbf{Formulation and scope.} Adapts an explicitly verified research-agenda passage. Mathematical objects, interventions, objectives and scope are proposed here; existing model-specific solutions are not relabeled unsolved.

\par\textbf{Open-question provenance.} An explicit question or research agenda was verified at the source locator. This does not by itself certify that every later version remains unresolved \sref{30007601}{1}.

\par\textbf{Historical priority.} This is the earliest explicit statement verified in this audit; first priority for the full catalogue formulation is not established.

\subsection{Short English statement}
In the explicitly specified setting: When does widespread AI trading improve information aggregation, and when does it create collusion, correlated errors or instability? The mathematical statement above fixes the target, assumptions and scope of this catalogue formulation.

\subsection{Sources}
\begin{itemize}\raggedright
\item[\textnormal{[S1]}]\hypertarget{source:30007601:1}{} Bank of England. 2026. Bank of England Agenda for Research, 2025–2028. The Financial System / Technological innovation, paragraphs 3–4; Prudential Architecture / Resilience, last paragraph.
\url{https://www.bankofengland.co.uk/research/bank-of-england-agenda-for-research}.
{\small\textit{Earliest verified question/agenda reference; historical priority qualified below. Requests algorithmic-market effects and AI resilience frameworks.}}
\end{itemize}
\end{document}
